\documentclass{article}
\usepackage[T1]{fontenc}
\usepackage{lmodern}
\usepackage{graphicx}
\usepackage{amsmath,amssymb}
\usepackage[a4paper,margin=2cm]{geometry}
\usepackage[absolute,overlay]{textpos}
\setlength{\TPHorizModule}{1mm}
\setlength{\TPVertModule}{1mm}
\usepackage[indonesian]{babel}
\usepackage{hyperref}
\setlength{\emergencystretch}{2em}
% unit: Halaman 8

\begin{document}

% === Halaman 8 (llm) ===

\section*{Prosedur Pembangkitan Sepasang Kunci}

\begin{enumerate}
  \item Pilih dua bilangan prima, $p$ dan $q$ (sebaiknya $p \neq q$)
  \item Hitung $n = pq$.
  \item Hitung $\phi(n) = (p-1)(q-1)$.
  \item Pilih sebuah bilangan bulat $e$ sebagai kunci publik, $e$ harus relatif prima terhadap $\phi(n)$.
  \item Hitung kunci dekripsi, $d$, dengan kekongruenan
  \[ ed \equiv 1 \pmod{\phi(n)} \text{ atau } d \equiv e^{-1} \pmod{\phi(n)} \to d = \frac{1+k\phi(n)}{e} \]
\end{enumerate}

Hasil dari algoritma di atas:
\begin{itemize}
  \item Kunci publik adalah pasangan $(e, n)$
  \item Kunci privat adalah pasangan $(d, n)$
\end{itemize}

\end{document}
